% Ad Familiares 13.21 --- headnote
% ad-familiares-13-21, 46 BC, Rome

\textit{Cicero to Servius Sulpicius Rufus, proconsul of
Achaia, written from Rome in 46 BC (the manuscript
dateline: \textit{Scr. Romae, ut videtur, a. 708 (46)}).
The second of the four Servius-recommendations preserved
in Book 13 (Fam. 13.20--23). The principals here are
two: M. Aemilius Avianius, a Roman businessman of long
acquaintance whose household and estate are at Sicyon
but who is for the moment detained at Cibyra in
Phrygia; and his freedman C. Avianius Hammonius, who is
managing his patron's affairs on the ground in Achaia.
The freedman is what the letter is really about ---
Cicero commends him in two capacities at once, as his
patron's procurator and on his own account.}

\textit{The interest of the piece is the personal claim
Cicero stakes on Hammonius's behalf: he stood by Cicero
"in my most distressing season" --- a transparent
reference to the exile of 58--57, or perhaps to the
civil-war years just past --- "with such fidelity and
good will, as if he had been freed by my own hand."
That last phrase is the rhetorical hinge: under Roman
patronage law, manumission created a lifelong bond of
duty between freedman and ex-master, and to say of
another man's freedman that he behaved \textit{as if} he
had been Cicero's own is the strongest form of
endorsement one Roman could offer another. The letter
also gives a clean small picture of how Roman commercial
houses operated in the Greek East: a Roman principal
based in Italy or temporarily abroad, a freedman
procurator on the spot, and a Roman provincial governor
whose goodwill smoothed everything.}
